\chapter{Agente de verificación FEVER}\label{cap:agente}

\section{Visión general}

El agente de verificación es la funcionalidad diferencial del plan
\emph{Super Pro}. Dado un texto libre, el sistema:

\begin{enumerate}[leftmargin=*,itemsep=0.2em]
  \item \textbf{Extrae \emph{claims}} verificables del texto.
  \item Para cada \emph{claim}, \textbf{recupera evidencias} mediante
        una API de búsqueda externa.
  \item Aplica un \textbf{clasificador NLI} sobre cada par
        (claim, evidencia) y obtiene
        $\hat{p}(\text{sup} \mid c, e_i)$, etc.
  \item \textbf{Agrega} las predicciones por evidencia en un
        veredicto por \emph{claim} con cuatro categorías:
        \textsc{supported}, \textsc{refuted}, \textsc{nei},
        \textsc{conflicting}.
  \item \textbf{Compone} un veredicto global a partir de los
        veredictos por \emph{claim}.
\end{enumerate}

\section{Componentes (módulo \texttt{fever})}

\begin{itemize}[leftmargin=*,itemsep=0.2em]
  \item \path{schemas.py}: \texttt{Claim}, \texttt{Evidence},
        \texttt{ClaimVerdict}, \texttt{VerificationReport}.
  \item \path{claim_extraction.py}: extracción de claims a partir del
        texto (heurística + LLM).
  \item \path{retrieval.py}: cliente de búsqueda (p.\ ej.\ Tavily) con
        cache opcional.
  \item \path{inference.py}: clases \texttt{FeverNLIClassifier},
        \texttt{OnnxNLIClassifier} y \texttt{StubNLIClassifier}, todas
        con método \texttt{predict(claim, evidence)}.
  \item \path{aggregation.py}: regla determinista de agregación con
        umbral configurable.
  \item \path{agent.py}: orquestador que ensambla los pasos.
\end{itemize}

\section{Decisiones de diseño}

\begin{itemize}[leftmargin=*,itemsep=0.2em]
  \item \textbf{Sin \emph{frameworks} de agentes}. Se evita LangChain
        para mantener trazabilidad de errores y simplicidad de
        despliegue.
  \item \textbf{Categoría \textsc{conflicting}}. Permite emitir un
        veredicto informativo cuando las evidencias se contradicen.
  \item \textbf{Inferencia desacoplada del modelo}. El clasificador
        NLI implementa una interfaz que admite tanto el modelo
        DeBERTa exportado a ONNX como un \emph{stub} determinista para
        pruebas.
      \item \textbf{Carga automática de ONNX}. En tiempo de arranque se
                        selecciona \texttt{OnnxNLIClassifier} cuando
                        \texttt{FEVER\_MODEL\_PATH} apunta a una carpeta con
                        \texttt{model.onnx}; en caso contrario se conserva el fallback
                        ligero para entornos de CI o desarrollo sin artefactos ML.
\end{itemize}

El modelo final integrado es \texttt{deberta-v3-small-fever-v1},
entrenado sobre FEVER-NLI y exportado a
\path{models/fever/deberta-v3-small-fever-v1}. En la partición de
prueba alcanza 0.7281 de exactitud y 0.7231 de F1 macro, frente a
0.5081 y 0.4842 del baseline TF-IDF + regresión logística. La prueba de
integración local cargó el ONNX con \texttt{onnxruntime} y produjo
veredictos coherentes para pares de ejemplo \textsc{sup} y
\textsc{ref}.

\section{Persistencia}

La migración \path{migrations/002_fever.sql} crea las tablas
\texttt{verification\_runs}, \texttt{claims} y \texttt{evidences} con
\emph{Row Level Security} restringido al \texttt{user\_id} del
solicitante.

\section{Tests}

17 tests Pytest cubren extracción, recuperación con \emph{stub},
inferencia, agregación y orquestación.
